\section{언어 코어와 합성 의미론}
\label{sec:semantics}

언어 코어는 상태를 가진 연산자와 스트림 연결로 프로그램의 의미를 구성한다.
Kahn의 수열 의미론은 통신선의 이력과 함수 합성을 연결하며,
동기 데이터흐름과 Lustre는 표본 발생 규칙과 논리 시간에 따라 연산자 그래프를 조직한다
\citep{kahn,sdf,lustre}.
본 절은 이 구조에 측정 스키마, 가용 시각, 데이터 계보를 결합한다.
형식 판단은 한 스키마 세대의 그래프에 적용하고, 설정 사건은 새로운 스키마 세대와
그에 대응하는 상태 초기화 규칙을 선택한다.

\subsection{원시 연산자와 상태 불변식}

\begin{definition}[상태 전이 연산자]
스트림 형식 $\tau$에 대해 $\Val(\tau)$를 해당 스키마와 시간 메타데이터를
만족하는 유한 블록의 집합으로 둔다. 원시 연산자 $o$는 다음 튜플을 가진다.
\begin{equation}
 o=(\tau_o^{\mathrm{in}},\tau_o^{\mathrm{out}},Z_o,z_o^0,
       \delta_o,\mathcal J_o,\mathcal E_o).
\end{equation}
$Z_o$는 상태 공간, $z_o^0$는 초기 상태,
$\mathcal J_o\subseteq Z_o$는 상태 불변식이다. 전이 함수는 입력 블록을 소비하고
다음 상태와 유한 개의 출력 블록을 계산한다.
\begin{equation}
 \delta_o:Z_o\times\Val(\tau_o^{\mathrm{in}})
 \longrightarrow Z_o\times\Val(\tau_o^{\mathrm{out}})^*.
\end{equation}
\end{definition}

각 원시 연산자의 계약은 초기 상태의 적합성과 전이의 폐쇄성을 포함한다.
\begin{align}
 z_o^0&\in\mathcal J_o,\label{eq:primitive-init}\\
 \left.
 \begin{aligned}
 z&\in\mathcal J_o,\qquad x\in\Val(\tau_o^{\mathrm{in}}),\\
 \delta_o(z,x)&=(z',Y)
 \end{aligned}\right\}
 &\Longrightarrow
 \begin{cases}
 z'\in\mathcal J_o,\\
 Y\in\Val(\tau_o^{\mathrm{out}})^*.
 \end{cases}
 \label{eq:primitive-preservation}
\end{align}
필터 상태는 지연 표본과 내부 누산값을 보유하고, 학습 연산자는 모델 매개변수와
갱신 계보를 보유한다. 스키마에 의존하는 상태는 채널 순서, 단위, 기준,
시계 세대의 식별자를 함께 저장한다.

출력 $y$의 실제 의존 레코드 집합을 $\Dep(y)$로 둔다.
연산자는 입력, 현재 상태, 매개변수에 연결된 조상 집합을 출력과 다음 상태로 전달한다.
이 계보 계약은 다음 포함 관계를 만족한다.
\begin{equation}
 \Dep(y)\subseteq\rho(y).
 \label{eq:lineage-inclusion}
\end{equation}
조건 분기와 매개변수 선택에 사용한 자료도 이 조상 집합에 포함한다.
따라서 계보는 수치 계산과 계산 경로 선택의 의존성을 함께 나타낸다.

\subsection{효과와 코어 문법}

\begin{definition}[연산 효과]
효과 $\mathcal E_o=(m_o,R_o,h_o)$는 처리 모드
$m_o\in\{\Live,\Record\}$, 요구 기능 집합 $R_o$,
출력의 측정 레이블을 기준으로 한 전방 관측 길이 $h_o$를 기록한다.
$\Live$ 모드는 가용 입력에 따라 결과를 생성하고, $\Record$ 모드는
입력 기록의 구간과 종료 상태를 연산 계약에 포함한다.
\end{definition}

유한 전방 창을 사용하는 연산자는 창에 필요한 표본과 완결 사건을 받은 뒤 출력을 생성한다.
창의 대표 측정시각은 끝점 anchor를 사용하며, 중심 창은 중심 anchor와
전방 관측 길이 $h_o$를 함께 명시한다.
출력 레코드는 대표 측정시각, 시간 지지 $\operatorname{Supp}(y)$,
실제 가용 시각을 함께 유지한다.
형식 검사기는 연산자의 시간 지원 구간을 합성하여 프로그램의 관측 범위를 계산한다.
연산 비용, 대기 시간, 상태 크기는 실행 계획의 자원 계약으로 연결한다.

프로그램 코어의 문법은 다음과 같다.
\begin{equation}
 \begin{aligned}
 p::={}&o\mid p_1;p_2\mid p_1\otimes p_2\\
       &\mid\operatorname{delay}_{d,z}
        \mid\operatorname{feedback}(p).
 \end{aligned}
 \label{eq:core-grammar}
\end{equation}
형식 환경 $\Gamma$는 원시 연산자의 서명과 변환 계약을 보관한다.
형식 판단 $\Gamma\vdash p:\tau\xrightarrow{\mathcal E}\sigma$는
프로그램 $p$의 입력 형식, 출력 형식, 실행 효과를 연결한다.
직렬 합성은 다음 규칙을 사용한다.
\begin{equation}
 \frac{
 \Gamma\vdash p:\tau_1\xrightarrow{\mathcal E_1}\tau_2
 \qquad
 \Gamma\vdash q:\tau_2\xrightarrow{\mathcal E_2}\tau_3
 }{
 \Gamma\vdash p;q:\tau_1
 \xrightarrow{\mathcal E_1\circ\mathcal E_2}\tau_3
 }.
 \label{eq:typing-sequential}
\end{equation}
공유 형식 $\tau_2$는 채널 순서, 측정량, 단위, 시계, 발생 규칙을 일치시킨다.
단위 변환, 재기준화, 재표본화, 시계 변환은 각각의 입력과 출력을 가진
명시적 연산자로 연결한다.
직렬 의미는 함수 합성으로 정의한다.
\begin{equation}
 \den{p;q}=\den{q}\circ\den{p}.
\end{equation}

병렬 합성은 입력과 출력 형식의 곱을 사용한다.
\begin{equation}
 \frac{
 \Gamma\vdash p:\tau_1\xrightarrow{\mathcal E_1}\sigma_1
 \qquad
 \Gamma\vdash q:\tau_2\xrightarrow{\mathcal E_2}\sigma_2
 }{
 \Gamma\vdash p\otimes q:\tau_1\times\tau_2
 \xrightarrow{\mathcal E_1\otimes\mathcal E_2}
 \sigma_1\times\sigma_2
 }.
 \label{eq:typing-parallel}
\end{equation}
두 합성에서 요구 기능은 $R_1\cup R_2$로 계산한다.
구성 연산자들이 $\Live$ 모드를 가지면 합성 결과에도 $\Live$를 부여한다.
$\Record$ 연산자를 포함하는 합성에는 $\Record$를 부여한다.
가변 상태와 장치 핸들은 지정한 actor가 소유하며, 읽기 전용 블록은
수명 계약에 따라 처리 분기에 공유한다. 인터페이스별 상태와 수명 계약은
RustBelt의 의미론적 라이브러리 검증 조건과 연결된다\citep{rustbelt}.

\subsection{가용 시각 인과성과 합성 보존}

동일한 초기 상태와 실행 환경에서 프로그램의 스트림 의미를 $\den{p}$로 쓴다.
각 유한 실행 시간 구간에 속하는 입력 레코드의 수는 유한하다.
실시간 연산자는 시각 $u$까지 가용한 자료와 실행 사건으로
해당 시각까지의 출력을 결정한다.

\begin{definition}[가용 시각 인과성]
스트림 변환 $F$가 모든 입력 $S$와 실행 시각 $u$에 대해
다음 절단 법칙을 만족하면 $F$를 가용 시각에 대해 인과적이라고 정의한다.
\begin{equation}
 \cut_u(F(S))=\cut_u\bigl(F(\cut_u(S))\bigr).
 \label{eq:availability-causality}
\end{equation}
각 출력의 가용 시각은 참조한 자료의 가용 시각과 다음 관계를 가진다.
\begin{equation}
 a_y\ge\max_{r\in\Dep(y)}a_r.
 \label{eq:availability-dependencies}
\end{equation}
상수 출력은 초기 상태와 상수 매개변수의 가용 시각을 사용한다.
\end{definition}

식~\eqref{eq:availability-causality}는 입력 관측과 출력 결정을 동일한
실행 시간 절단으로 연결한다. 측정 시각을 기준으로 한 창의 레이블은
창에 속한 물리 현상을 나타내고, 가용 시각은 그 창을 이용하는 의사결정의 시점을 나타낸다.
실행 기록은 이 두 시각과 모델 활성화 사건을 함께 보존한다.

\begin{proposition}[직렬ㆍ병렬 합성의 보존]
원시 연산자들이 식~\eqref{eq:primitive-init}--\eqref{eq:primitive-preservation}의
상태 계약과 식~\eqref{eq:lineage-inclusion}의 계보 계약을 만족한다고 하자.
원시 연산자, 직렬 합성, 병렬 합성으로 구성한 형식화된 프로그램은
유효 입력을 유효 출력으로 사상하고 출력의 실제 의존성을 계보에 포함한다.
구성 원시 연산자들이 식~\eqref{eq:availability-causality}를 만족하면
합성 프로그램도 가용 시각 인과성을 만족한다.
\end{proposition}

\begin{proof}
프로그램 구조에 대한 귀납법을 적용한다. 원시 연산자 경우는 해당 계약에서 따른다.
직렬 합성에서 첫 연산자의 출력 형식과 다음 연산자의 입력 형식이 일치하므로
유효 값의 사상이 이어진다. $F=\den{p}$, $G=\den{q}$라 두면
인과성은 다음 등식으로 얻는다.
\begin{align*}
 \cut_u(G(F(S)))
 &=\cut_u\bigl(G(\cut_u(F(S)))\bigr)\\
 &=\cut_u\bigl(G(\cut_u(F(\cut_u(S))))\bigr)\\
 &=\cut_u\bigl(G(F(\cut_u(S)))\bigr).
\end{align*}
계보의 합집합과 조상 폐쇄는 두 연산 단계의 실제 의존성을 포함한다.
병렬 합성은 두 성분에 귀납 가정을 각각 적용한다.
따라서 형식 적합성, 계보 포함 관계, 가용 시각 인과성을 전체 프로그램으로 확장한다.
\end{proof}

\subsection{초기 지연과 피드백 실행}

논리 틱으로 전개한 반응형 부분그래프는 틱 $k$의 유한 입력 블록과 actor 상태로
현재 출력과 다음 상태를 계산한다. 지연 연산자는 $d\ge1$개의 틱과
형식에 적합한 초기 블록 $z_0,\ldots,z_{d-1}$를 가진다.
\begin{equation}
 \operatorname{delay}_{d,z}x[k]
 =\begin{cases}
 z_k,&0\le k<d,\\
 x[k-d],&k\ge d.
 \end{cases}
 \label{eq:initialized-delay}
\end{equation}
입출력 시계와 스키마가 연결 형식을 결정하고, 초기 블록은 그 형식을 만족한다.
양의 지연을 가진 연결선을 제거한 그래프를 $G_0$라 둔다.
피드백 형식 규칙은 $G_0$의 비순환성과 각 actor의 결정적 전이를 요구한다.
이 규칙은 현재 틱의 계산 순서와 이전 틱에서 전달할 상태를 명시한다.
형식 $\tau\times\sigma\xrightarrow{\mathcal E}\nu\times\sigma$를 가진
부분그래프의 $\sigma$ 출력과 $\sigma$ 입력을 초기 지연으로 연결하면,
외부 형식은 $\tau\xrightarrow{\mathcal E}\nu$가 된다.

\begin{proposition}[지연 피드백의 실행 유일성]
반응형 그래프의 각 actor가 유효 입력과 현재 상태에 대해
전역적으로 정의된 결정적 전이를 가진다고 하자.
초기 상태와 지연 초기값이 각 형식을 만족하고 $G_0$가 비순환 그래프이면,
각 외부 입력 수열에 대해 유일한 형식화된 실행 수열이 존재한다.
\end{proposition}

\begin{proof}
틱 $k=0$의 지연값은 초기 블록으로 정해진다.
외부 입력과 지연값을 고정하고 $G_0$의 위상 순서대로 actor를 평가하면
모든 출력과 다음 상태가 유일하게 정해진다.
틱 $k$까지의 실행이 유일하게 정해졌다고 하자.
틱 $k+1$의 지연값은 초기 블록 또는 앞선 틱의 출력에서 결정된다.
같은 위상 순서 평가가 틱 $k+1$의 출력과 상태를 유일하게 결정한다.
각 actor의 상태 계약과 연결 형식은 유효 값을 다음 틱으로 전달한다.
자연수 귀납법으로 전체 실행 수열의 존재와 유일성을 얻는다.
\end{proof}

두 명제는 처리 그래프의 합성과 폐루프 실행에 대한 수학적 기준을 제공한다.
컴파일러는 형식화된 그래프를 실행 계획으로 옮기고,
그 계획에 토큰 수, 버퍼 용량, 호출 시간과 상태 소유권을 연결한다.
